\documentclass[12pt]{article}

% Layout.
\usepackage[top=1in, bottom=0.75in, left=0.75in, right=0.75in, headheight=1in, headsep=6pt]{geometry}

% Fonts.
\usepackage{mathptmx}

\usepackage[scaled=0.86]{helvet}
\renewcommand{\emph}[1]{\textsf{\textbf{#1}}}

% Misc packages.
\usepackage{amsmath,amssymb,latexsym,amsthm}
\usepackage{graphicx,hyperref}
\usepackage{array}
\usepackage{xcolor}
\usepackage{multicol}
\usepackage{tabularx,colortbl}
\usepackage{enumitem}
\usepackage{soul,multicol}
\usepackage{setspace}
\doublespacing

\hypersetup{
    colorlinks=true,
    linkcolor=blue,
    filecolor=magenta,      
    urlcolor=blue,
    pdftitle={Proofs Worksheet},
    pdfpagemode=FullScreen,
    }

\def\mailto#1{\href{mailto:#1}{#1}}

%% special math symbols
\newcommand{\N}{\mathbb{N}}
\newcommand{\Z}{\mathbb{Z}}
\newcommand{\R}{\mathbb{R}}
\newcommand{\Q}{\mathbb{Q}}
\newcommand{\sse}{\subseteq}
\newcommand{\mcp}{\mathcal{P}}
\newcommand{\ol}[1]{\overline{#1}}

%other specials
\newcommand{\be}{\begin{enumerate}}
\newcommand{\ee}{\end{enumerate}}
\newcommand{\se}{\subseteq}
\newcommand{\spe}{\supseteq}
\newcommand{\ds}{\displaystyle}
\newcommand{\lcm}{\text{lcm}}
%\newcommand{\gcd}{\text{gcd}}

% Paragraph spacing
\parindent 0pt
\parskip 6pt plus 1pt
\def\tableindent{\hskip 0.5 in}
\def\ts{\hskip 1.5 em}

%header
\usepackage{fancyhdr}
\pagestyle{fancy} 
\lhead{\large\sf\textbf{MATH 265: Introduction to Mathematical Proofs}}
\rhead{\large\sf\textbf{Worksheet 19: Ch 11}}

\newcommand{\localhead}[1]{\par\smallskip\textbf{#1}\nobreak\\}%
\def\heading#1{\localhead{\large\emph{#1}}}
\def\subheading#1{\localhead{\emph{#1}}}

\newenvironment{clist}%
{\bgroup\parskip 0pt\begin{list}{$\bullet$}{\partopsep 4pt\topsep 0pt\itemsep -2pt}}%
{\end{list}\egroup}%

\begin{document}
\begin{enumerate}
\item A \emph{relation} $R$ on $A$ is 
\vfill

\item Examples
\vfill

\vfill

\item Suppose $R$ is a relation on the set $A.$
	\begin{enumerate}
	\item We say $R$ is \emph{symmetric} if 
	\vfill

	\item We say $R$ is \emph{reflexive} if 
	\vfill

	\item We say $R$ is \emph{transitive} if 
	\vfill
	\end{enumerate}
\newpage
\item Let $R$ be the relation on $A=\{0,1,2,3,4\}$ defined as $a\: R \: b$ if $a \mid b.$ List all the elements in $R$. What is the smallest positive number you could add to the set $A$ that would increase the order $R$ as much as possible?
\vfill
\item Let $R$ be a relation on $A=\Z$ defined by $a \;R\; b$ if $a \equiv b \pmod 3.$ Find three elements of $R$ that contain zero in the first coordinate, find three elements of $R$ that contain 1 in the second coordinate, and find three ordered pairs in $\Z \times \Z$ that are not in $R.$
\vfill
\item Suppose $R$ is a relation on $\mathcal{P}(A)$ where $A=\{0,1,2,3,4\}$ defined as $X \: R \: Y$ if $X \cap Y \not = \emptyset.$ List 5 elements in $R$ and 3 elements of $ \mathcal{P}(A)\times \mathcal{P}(A)$ that are not in $R.$
\vfill
\item For each relation above determine if it is symmetric, reflexive or transitive. Make a formal argument/proof of your conclusion.
\end{enumerate}
\end{document}










